\section{算力分配：1T 入场券与 3:1:1}

访谈中最具体的资源判断之一，是“1T 模型是入场券”和算力比例变化。罗福莉认为，如果要做到接近 Claude Opus 4.6 的 Agent 水平，1T 级基础模型可能是入场券。但进入 Agent 阶段后，算力分配不能仍然由预训练独占。

\begin{figure}[H]
\centering
\includegraphics[width=0.92\textwidth]{compute-allocation.png}
\caption{从 Chat 阶段到 Agent 阶段的算力分配趋势。}
\end{figure}

\begin{knowledgebox}{读图：3:5:1 与 3:1:1 的含义}
图中 3:5:1 和 3:1:1 是访谈中的口头比例，用来表达资源重心变化。Chat 时代预训练占比更高；Agent 时代 research、pre-train、post-train 的关系更接近均衡，后训练用卡权重显著上升。
\end{knowledgebox}

\subsection{1T 是什么意义上的入场券}

“1T 模型是入场券”不是说只要 1T 参数就够，也不是说小模型没有价值。它指向一个前沿 Agent 能力门槛：若要对标最强 Agent 系统，基座模型规模、知识、推理、代码和多步任务能力必须达到足够高的底座水平。后训练和框架可以放大模型，但不能凭空制造所有能力。

\begin{warningbox}{规模不是唯一答案}
参数量只是入口变量。Agent 系统还受训练数据、上下文、代码能力、工具使用、后训练、任务环境、成本和组织执行影响。把“1T 入场券”理解成“堆参数即可领先”是误读。
\end{warningbox}

\subsection{Post-train 与 pre-train 的平衡}

罗福莉认为，顶尖团队预训练与后训练的资源比例已经接近 1:1。这意味着后训练不再是预训练后的尾声，而是与预训练并列的主战场。尤其在 Agent 任务中，后训练直接决定模型是否能适配环境、工具和长程任务。

\begin{importantbox}{资源配置是战略判断}
卡怎么分，不是财务问题，而是技术路线判断。把卡继续全部堆给预训练，和把足够卡留给 research、post-train、rollout、eval，会得到完全不同的组织能力。
\end{importantbox}

\subsection{本章小结}

Agent 时代的算力分配更均衡、更复杂。前沿竞争既需要大基座，也需要后训练和研究资源。3:1:1 的核心信息是：后训练已经从配角变成主角。

